%================================
% Vocabulary & Language Focus — XB2 (Unit 2)
% 编译方式: xelatex main.tex (需运行两遍以生成目录)
%================================
\documentclass[../main.tex]{subfiles}

\begin{document}

\section{Unit 2}

\vw{solidarity}{/ˌsɒlɪˈdærəti/}{n.}

\sense{n.}{团结，齐心协力}
\col{solidarity with sb}
\ex{The workers showed solidarity with the strikers.}
\col{express/declare solidarity}
\ex{Thousands marched to express solidarity with the victims.}
\col{in solidarity with sb}
\ex{He resigned in solidarity with his colleague.}

\begin{origin}
来自法语 \textit{solidarité}，源于形容词 \textit{solidaire}（连带责任的、团结的），再追溯到拉丁语 \textit{solidus}（坚固的、实心的）。同源词：\textbf{solid}（固体的、坚实的）、\textbf{consolidate}（巩固）、\textbf{solder}（焊接）。"团结"就是"抱成一个坚实的整体"。
\end{origin}

\vw{participate}{/pɑːˈtɪsɪpeɪt/}{vi.}

\sense{vi.}{参加，参与}
\col{participate in sth}
\ex{All students are encouraged to participate in the competition.}
\col{participate actively/fully}
\ex{She participated actively in the discussion.}
\col{participation (n.) 参与}
\ex{participation in class activities}

\col{participate 是不及物动词，必须接 in}

\begin{origin}
来自拉丁语 \textit{participare}（分享、分得一份）：\textit{pars}（部分）+ \textit{capere}（取）。同源词：\textbf{part}、\textbf{particle}（微粒）、\textbf{participation}。"取一份、出一份力"就是参与。
\end{origin}

\begin{lex}{participate / take part in / attend / join}
\textbf{participate in} 与 \textbf{take part in} 同义，强调"起作用、出力气"；\textbf{attend} 只表"到场"（\textit{attend a lecture} 听讲座，未必发言）；\textbf{join} 表"加入组织/人群"（\textit{join a club}；\textit{join us}）。"参加比赛"用 \textit{take part in a contest} 或 \textit{compete in}。
\end{lex}

\vw{racial}{/ˈreɪʃl/}{adj.}

\sense{adj.}{种族的；种族间的}
\col{racial discrimination/prejudice 种族歧视；种族偏见}
\ex{The government promised to fight racial discrimination.}
\col{racial equality/harmony 种族平等；种族和睦}
\ex{The school celebrates racial harmony among its students.}
\col{racial tension 种族关系紧张}
\ex{Racial tension flared up after the incident.}

\begin{origin}
race（人种）的词源有争议，一说借自意大利语 \textit{razza}（种类、血统），或与拉丁语 \textit{ratio}（类别、性质）有关；racial 即 race 的形容词（-ial 后缀，同 legal、social）。"人种"概念的种族色彩是近代欧洲学术语境强加的。
\end{origin}

\vw{diverse}{/daɪˈvɜːs/}{adj.}

\sense{adj.}{多种多样的，不同的}
\col{a diverse range of sth}
\ex{The city has a diverse range of cultures and cuisines.}
\col{diverse opinions/backgrounds}
\ex{Students come from diverse cultural backgrounds.}
\col{diversity (n.) 多样性}
\ex{cultural diversity}

\begin{origin}
来自拉丁语 \textit{diversus}（朝不同方向的），是 \textit{divertere}（转向不同方向）的过去分词：\textit{dis-}（分开）+ \textit{vertere}（转）。同源词：\textbf{divert}（使转向）、\textbf{divorce}（离婚：把两个人转向两边）、\textbf{diversity}。"各转各的方向"自然"多种多样"。
\end{origin}

\vw{joint}{/dʒɔɪnt/}{adj. \& n.}

\sense{adj.}{联合的，共同的}
\col{a joint venture 合资企业}
\ex{The two firms set up a joint venture in Shanghai.}
\col{a joint statement/declaration 联合声明}
\ex{The two countries issued a joint statement.}
\col{joint efforts 共同努力}
\ex{The project was completed through our joint efforts.}

\sense{n.}{关节}
\col{joint pain/inflammation 关节疼痛；关节炎}
\ex{He suffers from joint pain in cold weather.}

\sense{n.}{接合处}
\ex{The pipes are leaking at the joints.}

\begin{origin}
来自古法语 \textit{jointe}（接合处），源于 \textit{joindre}（连接），再追溯到拉丁语 \textit{iungere}（连接）。同源词：\textbf{join}、\textbf{junction}（交叉口）、\textbf{conjunction}（连接词）、\textbf{conjugal}（婚姻的，共同驾驭的）。"连接之处"既是关节，也是共同的事业。
\end{origin}

\vw{motivate}{/ˈməʊtɪveɪt/}{vt.}

\sense{vt.}{激励，激发}
\col{motivate sb to do sth}
\ex{Good teachers know how to motivate students to learn.}
\col{motivated (adj.) 有干劲的}
\ex{The team is highly motivated.}

\sense{vt.}{成为……的动机，是……的原因}
\col{be motivated by sth}
\ex{He was motivated by a desire for wealth.}
\col{motivation (n.) 动机；积极性}
\ex{What is your motivation for learning English?}

\begin{origin}
来自拉丁语 \textit{motivus}（运动的），源于 \textit{movere}（移动）。同源词：\textbf{move}、\textbf{motion}（运动）、\textbf{emotion}（情绪：e- 向外 + motion $\to$ 涌动的心理活动）、\textbf{motor}。"推动某人行动"即激励。
\end{origin}

\vw{demonstrate}{/ˈdemənstreɪt/}{vt. \& vi.}

\sense{vt.}{表现，表露}
\col{demonstrate one's support/skill}
\ex{She demonstrated great courage during the crisis.}

\sense{vt.}{证实，证明}
\col{demonstrate that ...}
\ex{The study demonstrates that the drug is effective.}
\ex{These figures demonstrate the success of the campaign.}

\sense{vt.}{示范，演示}
\col{demonstrate how to do sth}
\ex{Let me demonstrate how to use the machine.}

\sense{vi.}{游行示威}
\col{demonstrate against sth}
\ex{Thousands demonstrated against the new law.}

\begin{origin}
来自拉丁语 \textit{demonstrare}（指出、表明）：\textit{de-}（完全）+ \textit{monstrare}（显示），词根 \textit{monstrum}（预兆、怪异的显现）。同源词：\textbf{monster}（怪物——神示的异象）。"把东西明明白白指给人看"——演示、证明、示威都共享这个意象。
\end{origin}

\vw{afterwards}{/ˈɑːftəwədz/}{adv.}

\sense{adv.}{以后，后来}
\col{soon/shortly afterwards}
\ex{He arrived shortly afterwards.}
\ex{We had dinner, and afterwards we went for a walk.}

\begin{gram}[afterwards 与 after]
\textit{afterwards} 是副词，单独使用（\textit{We met afterwards.}），\textcolor{WildStrawberry!80!black}{\ding{55}} \textit{afterwards the meeting}；\textit{after} 是介词或连词，须接宾语或从句（\textit{after the meeting}；\textit{after the meeting ended}）。美式英语常写作 \textit{afterward}。同类构词：\textit{afterward(s)} = after + -ward(s)（朝向）。
\end{gram}

\vw{proceed}{/prəˈsiːd/}{vi.}

\sense{vi.}{接着做，继而做}
\col{proceed to do sth}
\ex{After the applause died down, she proceeded to read her speech.}

\sense{vi.}{继续从事（或进行）}
\col{proceed with sth}
\ex{We decided to proceed with the plan despite the risks.}

\sense{vi.}{行进，前往}
\col{proceed to + 地点}
\ex{The tour group proceeded to the next town.}

\begin{gram}[proceed to do 与 continue doing]
\textit{proceed to do sth} 强调"接着去做另一件事"（\textit{He then proceeded to explain}），暗含"完成上一步、进入下一步"；\textit{continue doing sth} 强调"不间断地继续同一件事"。名词形式：\textbf{process}（过程）、\textbf{procedure}（程序）、\textbf{proceedings}（议程）。
\end{gram}

\begin{origin}
来自拉丁语 \textit{procedere}（向前走）：\textit{pro-}（向前）+ \textit{cedere}（走、让步）。同源词：\textbf{process}、\textbf{precede}（pre- 先 + cede $\to$ 先走 $\to$ 领先）、\textbf{recede}（后退）、\textbf{succeed}（suc- 紧跟 + cede $\to$ 跟着走 $\to$ 接替、成功）。\textit{cedere}（行走）是英语中最能"走路"的词根之一。
\end{origin}

\vw{flame}{/fleɪm/}{n.}

\sense{n.}{火焰，火舌}
\col{burst into flames 突然起火}
\ex{The plane burst into flames after the crash.}
\col{in flames 燃烧着}
\ex{The whole building was in flames.}
\col{the flames of sth（比喻）激情的烈焰}
\ex{the flames of passion}

\col{an old flame 旧情人（口语）}
\ex{He met an old flame at the reunion.}

\begin{origin}
来自拉丁语 \textit{flamma}（火焰），源于 \textit{flagrare}（燃烧），词根可追溯到原始印欧语 *\textit{bhel-}（发光、燃烧）。同源词：\textbf{flagrant}（罪恶昭彰的：烧到明面上的）、\textbf{conflagration}（大火）、\textbf{flammable}（易燃的）。同一词根还给了英语 \textbf{blond}（金发的，闪亮的）和 \textbf{blue}。
\end{origin}

\vw{applaud}{/əˈplɔːd/}{vt. \& vi.}

\sense{vt. \& vi.}{称赞，赞许；鼓掌}
\col{applaud sb/sth（鼓掌）}
\ex{The audience applauded the performers warmly.}
\col{applaud a decision/effort（称赞）}
\ex{We applaud the government's decision to invest in education.}
\col{applause (n.) 掌声}
\ex{a round of applause}

\begin{origin}
来自拉丁语 \textit{applaudere}（拍手称快）：\textit{ad-}（向）+ \textit{plaudere}（拍手）。同源词：\textbf{applause}、\textbf{plaudits}（喝彩）、\textbf{explode}（ex- 向外 + plaudere 拍手 $\to$ 把演员"拍"下台 $\to$ 爆发）。"鼓掌"与"爆炸"竟是同根——喝倒彩的掌声轰走演员，explode 由此而来。
\end{origin}

\vp{track and field}

\sense{n.}{田径运动}
\col{track and field events 田径项目}
\ex{She excels at track and field.}
\ex{He won three gold medals in track and field events.}

\begin{origin}
track（跑道）+ field（场地）：田径分为径赛（track events，在跑道上赛跑）和田赛（field events，在场地里跳跃、投掷）。汉语"田径"正是 track and field 的直译。
\end{origin}

\vw{consistent}{/kənˈsɪstənt/}{adj.}

\sense{adj.}{一致的，始终如一的}
\col{be consistent with sth 与……一致}
\ex{The results are consistent with our expectations.}
\col{consistent effort/performance 坚持不懈的}
\ex{She has been consistent in her efforts.}

\sense{adj.}{连续的}
\col{consistent rain/attacks}
\ex{Consistent rain caused flooding in the lowlands.}

\col{consistency (n.) 一致性}

\begin{lex}{consistent / constant}
\textbf{consistent} 强调"前后一致、不自相矛盾"（\textit{consistent with the facts}）；\textbf{constant} 强调"持续不断、从不停止"（\textit{constant rain} 雨下个不停；\textit{a constant companion} 不离不弃的伴侣）。
\end{lex}

\begin{origin}
来自拉丁语 \textit{consistere}（站在一起、站得住）：\textit{con-}（一起）+ \textit{sistere}（站立）。同源词：\textbf{consist}（由……组成）、\textbf{resist}（反抗：站在对面）。"站在一起不散架"自然"始终如一"。
\end{origin}

\vp{call on sb to do sth}

\sense{v.}{正式邀请某人做某事；请求某人做某事}
\ex{The UN called on both sides to end the fighting.}
\ex{The president called on Congress to pass the bill.}

\col{call on/upon sb（还可表"拜访某人"）}
\ex{We called on our grandparents on the way home.}

\begin{origin}
call on 由"上门拜访"引申出"郑重请求"——登门相请，比 ask 更有分量。注意 \textbf{call at} 接地点（\textit{call at a shop}），\textbf{call on} 接人。
\end{origin}

\vw{humanity}{/hjuːˈmænəti/}{n.}

\sense{n.}{（统称）人，人类}
\ex{The vaccine is a gift to all humanity.}
\col{a crime against humanity 反人类罪}

\sense{n.}{人性}
\ex{War brings out the worst in humanity.}

\sense{n.}{人道，仁慈}
\col{treat sb with humanity}
\ex{The prisoners were treated with humanity.}

\col{the humanities 人文学科}
\ex{She studied the humanities at university.}

\begin{origin}
来自拉丁语 \textit{humanitas}（人性、教养），源于 \textit{humanus}（人的），词根 \textit{homo}（人）。同源词：\textbf{human}、\textbf{humane}（人道的）、\textbf{homicide}（homo 人 + cide 杀 $\to$ 杀人）。罗马人把"人之所以为人"理解为教养与文雅，humanities（人文学科）由此得名。
\end{origin}

\vw{rank}{/ræŋk/}{vi. \& vt. \& n.}

\sense{vi.}{属于某等级}
\col{rank among/above}
\ex{The university ranks among the top ten in the country.}
\col{rank as}
\ex{He ranks as one of the greatest players of all time.}

\sense{vt.}{把……分等级}
\col{rank sb/sth as sth}
\ex{The magazine ranks the city as the best place to live.}

\sense{vt.}{排列，使排成行}
\ex{The soldiers were ranked in neat rows.}

\sense{n.}{地位，级别；等级；军衔}
\col{rise through the ranks 逐级晋升}
\ex{He rose through the ranks to become a general.}
\col{the rank of sth}
\ex{the rank of colonel}

\sense{n.}{排，列}
\ex{a rank of soldiers}

\begin{origin}
来自古法语 \textit{ranc}（行列），源于法兰克语（日耳曼语）*\textit{hring}（圆圈、队伍）——与 \textbf{ring}（圈）同源。同源词：\textbf{arrange}（a- + rengier 排成行 $\to$ 安排）。"排成一圈、排成一行"→ 行列 → 等级。
\end{origin}

\vw{trend}{/trend/}{n.}

\sense{n.}{趋势，趋向}
\col{a trend towards sth}
\ex{There is a trend towards remote working.}
\col{an upward/downward trend}
\ex{The stock market shows an upward trend.}
\col{set a trend 引领潮流}
\ex{The company set a trend in smartphone design.}
\col{follow the trend 跟风}

\begin{origin}
来自中古英语 \textit{trenden}（旋转、转动），可追溯到原始印欧语 *\textit{trend-}（转动）。由"转动的方向"引申出"动向、趋势"，19 世纪又发展出"时尚潮流"之义。
\end{origin}

\vw{hurdle}{/ˈhɜːdl/}{n.}

\sense{n.}{(hurdles) 跨栏赛；栏架}
\col{the 400-metre hurdles}
\ex{She won gold in the 400-metre hurdles.}
\col{clear a hurdle 跨过栏架}

\sense{n.}{障碍}
\col{overcome a hurdle}
\ex{Funding remains the biggest hurdle to the project.}

\begin{origin}
源自古英语 \textit{hyrdel}（用枝条编成的临时栅栏），词根有"编织"之意（原始印欧语 *\textit{kert-} 编织）。临时栅栏本用于圈牲口，后被跨栏比赛借用——"跨过栅栏"成了"跨过障碍"。
\end{origin}

\vw{election}{/ɪˈlekʃn/}{n.}

\sense{n.}{选举，推选；当选}
\col{hold an election 举行选举}
\ex{The election will be held in May.}
\col{win/lose an election}
\ex{She won the election by a narrow margin.}
\col{run in an election 参加竞选}
\ex{He ran in the local election last year.}

\begin{origin}
来自拉丁语 \textit{electio}（挑选），源于 \textit{eligere}（选出）：\textit{e-}（向外）+ \textit{legere}（选择、采集、阅读）。同源词：\textbf{elect}（选举）、\textbf{select}（选择）、\textbf{collect}（收集）、\textbf{lecture}（lecture：legere 读 $\to$ 采集知识 $\to$ 讲授）。"选出来的人"就是当选者。
\end{origin}

\vw{advent}{/ˈædvent/}{n.}

\sense{n.}{（重要事件、人物等的）到来，出现}
\col{the advent of sth}
\ex{The advent of the Internet changed the world.}
\ex{the advent of spring}

\col{Advent 基督降临节（宗教）}

\begin{origin}
来自拉丁语 \textit{adventus}（到来），源于 \textit{advenire}（来到）：\textit{ad-}（向）+ \textit{venire}（来）。同源词：\textbf{venue}（人们到来的地方）、\textbf{convenient}（con- 一起 + venire 来 $\to$ 凑到一起合适 $\to$ 方便的）、\textbf{event}（e- 出 + venire $\to$ 走出来的事 $\to$ 事件）、\textbf{adventure}（adventurus 即将到来的 $\to$ 奇遇）。\textit{venire}（来）是拉丁语又一支多产词根。
\end{origin}

\vw{faith}{/feɪθ/}{n.}

\sense{n.}{信心}
\col{have faith in sb/sth}
\ex{I have great faith in your ability.}
\col{lose faith in sth}
\ex{People are losing faith in the system.}

\sense{n.}{宗教信仰}
\col{a man/woman of faith 有信仰的人}
\ex{the Christian faith}

\begin{origin}
来自古法语 \textit{feid}，源于拉丁语 \textit{fides}（信任、忠诚），词根 \textit{fidere}（相信）。同源词：\textbf{fidelity}（忠贞）、\textbf{confident}（con- 完全 + fidens 相信 $\to$ 自信）、\textbf{confide}（吐露心声）。原始印欧语词根 *\textit{bheidh-} 意为"信任"。
\end{origin}

\vw{ambition}{/æmˈbɪʃn/}{n.}

\sense{n.}{追求的目标，夙愿}
\col{an ambition to do sth}
\ex{Her ambition is to become a doctor.}
\col{achieve/realize one's ambition}
\ex{He finally achieved his lifelong ambition.}

\sense{n.}{野心，雄心，壮志}
\col{a man of ambition}
\ex{Ambition alone is not enough without talent.}

\col{ambitious (adj.) 有雄心的；野心勃勃的}

\begin{origin}
来自拉丁语 \textit{ambitio}，本义是"四处奔走"——\textit{ambire}（绕行）：\textit{ambi-}（周围）+ \textit{ire}（走）。古罗马竞选公职的人四处奔走拉票，ambitio 遂由"到处跑"转指"野心、抱负"。同源词：\textbf{ambient}（环绕的）。
\end{origin}

\vw{ascend}{/əˈsend/}{vi. \& vt.}

\sense{vi. \& vt.}{上升，升高，登高}
\col{ascend a mountain/stairs}
\ex{They ascended the mountain before dawn.}
\col{ascend the throne 登基}
\ex{She ascended the throne at the age of eighteen.}
\ex{The plane ascended rapidly into the clouds.}
\col{in ascending order 按升序}

\begin{lex}{ascend / climb / mount}
\textbf{ascend} 最正式、书面（\textit{ascend the throne}）；\textbf{climb} 日常通用，常指手脚并用地爬（\textit{climb a tree}）；\textbf{mount} 侧重"登上"（\textit{mount a horse} 上马）。反义词：descend（下降）、dismount（下马）。
\end{lex}

\begin{origin}
来自拉丁语 \textit{ascendere}（攀上）：\textit{ad-}（向）+ \textit{scandere}（攀登）。同源词：\textbf{descend}（下降）、\textbf{scale}（拉丁语 scala 梯子，源于 scandere——"攀登的器具"）。梯子（scale）与攀登（ascend）同出一根。
\end{origin}

\vw{identical}{/aɪˈdentɪkl/}{adj.}

\sense{adj.}{完全同样的}
\col{be identical to/with sth}
\ex{Her dress is identical to mine.}
\col{identical twins 同卵双胞胎}
\ex{The identical twins are hard to tell apart.}

\begin{origin}
来自拉丁语 \textit{identitas}（同一性），源于代词 \textit{idem}（同一个）。同源词：\textbf{identity}（身份）、\textbf{identify}（辨认）。"同一个"就是 identical——一模一样。
\end{origin}

\vw{fellow}{/ˈfeləʊ/}{adj. \& n.}

\sense{adj.}{同伴的，同情况的，同类的}
\col{fellow students/workers/citizens 同学；工友；同胞}
\ex{She discussed the idea with her fellow students.}

\sense{n.}{同辈，同类；男人，家伙}
\ex{He's a friendly fellow.}

\col{fellowship (n.) 交情；团体}

\begin{origin}
源自古英语 \textit{feolaga}（合伙人），借自古诺斯语 \textit{félagi}：\textit{fé}（财产、钱）+ \textit{lag}（放置）——字面义"把财产放在一起的人"。同源词：\textbf{fee}（费用）。"共同出资的人"就是同伴。
\end{origin}

\vw{relay}{/ˈriːleɪ/ \ n. \ \ /ˈriːleɪ; rɪˈleɪ/ \ vt.}{n. \& vt.}

\sense{n.}{接力赛}
\col{a relay race}
\ex{She ran the last leg of the relay.}
\col{in relays 轮班，分批}
\ex{The volunteers worked in relays through the night.}

\sense{vt.}{转发（信息、消息等）；播放，转播}
\col{relay a message to sb}
\ex{Could you relay the message to the manager?}
\col{relay the match live 现场转播}

\begin{origin}
来自古法语 \textit{relais}（换马的驿站），源于 \textit{relayer}（换马）：\textit{re-}（再）+ \textit{laier}（留下）。古代驿路上每隔一段设驿站换马接力，信息由此一站站传下去——"接力赛"和"转播"都是这个意象。
\end{origin}

\vw{chest}{/tʃest/}{n.}

\sense{n.}{胸部}
\col{chest pain 胸痛}
\ex{He complained of chest pain after exercise.}

\sense{n.}{大箱子}
\col{a chest of drawers 五斗柜}
\ex{The documents were kept in a wooden chest.}

\col{get sth off one's chest 一吐为快}
\ex{He felt better after getting it off his chest.}

\begin{origin}
源自古英语 \textit{cest}，借自古法语 \textit{chiste}，源于拉丁语 \textit{cista}（箱子、篮子），再追溯到希腊语 \textit{kistē}（篮子）。"胸膛"像一只收拢的箱子——两个义项同出一源。
\end{origin}

\vw{amateur}{/ˈæmətə(r)/}{n. \& adj.}

\sense{n.}{业余运动员，业余爱好者；生手，外行}
\ex{The tournament is open to both amateurs and professionals.}
\ex{When it comes to car repair, I'm a complete amateur.}

\sense{adj.}{业余的}
\col{an amateur photographer/boxer}

\col{amateurish (adj.) 业余水平的，外行的}

\begin{origin}
借自法语 \textit{amateur}（爱好者），源于拉丁语 \textit{amator}（热爱者），词根 \textit{amare}（爱）。同源词：\textbf{amour}（爱情）、\textbf{amiable}（和蔼可亲的）、\textbf{amicable}（友好的）。"出于热爱而从事"的人就是业余爱好者——与"为钱而做"的 professional（职业选手）相对。
\end{origin}

\vw{gradually}{/ˈɡrædʒuəli/}{adv.}

\sense{adv.}{逐渐地，逐步地}
\ex{The weather gradually improved.}
\ex{She gradually got used to the new job.}
\col{gradual (adj.) 逐渐的}
\ex{a gradual change in temperature}

\begin{origin}
来自拉丁语 \textit{gradus}（步）$\to$ \textit{gradualis}（一步一步的）$\to$ gradually。同源词：\textbf{grade}（年级、等级）、\textbf{progress}（pro- 前 + gress 步 $\to$ 前进）、\textbf{graduate}（毕业：走完学业的台阶）、\textbf{aggressive}（ad- 向 + gress 步 $\to$ 步步进逼 $\to$ 好斗的）。"一步一个脚印"就是 gradually。
\end{origin}

\vw{session}{/ˈseʃn/}{n.}

\sense{n.}{一段时间，一场}
\col{a training session 训练课}
\ex{We had a two-hour training session this morning.}
\col{a photo/recording session 拍照；录音时段}

\sense{n.}{会议，开庭}
\col{in session 在开会；在开庭}
\ex{The court is now in session.}
\col{a session of Parliament 议会会议}

\sense{n.}{学年}
\col{the summer session 夏季学期}

\begin{origin}
来自拉丁语 \textit{sessio}（坐着），源于 \textit{sedere}（坐）。同源词：\textbf{sit}、\textbf{seat}、\textbf{president}（preside 坐在前面主持 $\to$ 总统、主席）。"坐着开完的一段时间"就是 session。
\end{origin}

\vw{limit}{/ˈlɪmɪt/}{n. \& vt.}

\sense{n.}{限度，限制}
\col{the speed limit 限速}
\ex{He was fined for exceeding the speed limit.}
\col{within limits 在限度内}

\sense{n.}{限量，限额}
\col{a time limit 时限}
\ex{There is a time limit of 30 minutes for the test.}

\sense{n.}{界限}
\col{the limits of sth}
\ex{There are limits to what we can achieve.}

\sense{vt.}{限制，限定}
\col{limit sth to sth}
\ex{Please limit your answer to 200 words.}

\sense{vt.}{限量，减量}
\ex{Doctors advise limiting salt intake.}

\begin{origin}
来自拉丁语 \textit{limes}（田界、边界；复数 \textit{limites}）。古罗马帝国的边境界墙就叫 limes（莱姆斯防线）。同源词：\textbf{limit}、\textbf{delimit}（划界）。"界限"是罗马人丈量帝国土地的术语。
\end{origin}

\vw{constitution}{/ˌkɒnstɪˈtjuːʃn/}{n.}

\sense{n.}{体质，身体素质}
\col{a strong/weak constitution}
\ex{He has a strong constitution and rarely falls ill.}

\sense{n.}{宪法，章程}
\col{the Constitution（某国）宪法}
\ex{The Constitution guarantees freedom of speech.}
\col{the constitution of a club 社团章程}

\begin{origin}
来自拉丁语 \textit{constitutio}（制度、规定），源于 \textit{constituere}（确立）：\textit{con-}（一起）+ \textit{statuere}（设立），词根 \textit{stare}（站立）。同源词：\textbf{statute}（法令）、\textbf{institute}（学院；制定）、\textbf{constitute}（构成）。"立起来的制度"就是宪法；"立起来的身体"就是体质。
\end{origin}

\vw{pour}{/pɔː(r)/}{vt. \& vi.}

\sense{vt.}{倒，斟}
\col{pour sth into/out of sth}
\ex{She poured the tea into the cups.}
\col{pour sb sth}
\ex{Let me pour you a drink.}

\sense{vt.}{倾倒，倒出}
\ex{He poured the milk down the sink.}

\sense{vi.}{涌流，倾泻}
\col{pour in/out/down}
\ex{Thousands of letters poured in from all over the country.}
\ex{Sweat poured down his face.}

\sense{vi.}{下大雨}
\col{It is pouring (with rain).}
\ex{It poured with rain all day.}
\col{the rain poured down}

\begin{gram}[it 作主语表天气]
表"下大雨"须用 it 作形式主语：\textit{It's pouring (down).}（下着倾盆大雨）；名词形式 \textbf{downpour}（暴雨）：\textit{a sudden downpour}。这类"it 指天气"的句型中 it 无语义，只充当语法主语，同类还有 \textit{It's raining cats and dogs.}
\end{gram}

\begin{origin}
来自古法语 \textit{purer}（过滤、倒出），源于拉丁语 \textit{purare}（使纯净），词根 \textit{purus}（纯净的）。同源词：\textbf{pure}、\textbf{purify}（净化）。13 世纪由"滤清液体"转义为"倾倒"。
\end{origin}

\vp{walk on air}

\sense{v.}{欢天喜地，得意扬扬}
\ex{She was walking on air after winning the scholarship.}

\begin{origin}
字面义"踩着空气走路"——高兴得脚不沾地。同类"高兴得飘起来"的习语：\textit{on cloud nine}（在九霄云外）。
\end{origin}

\vp{keep sth in perspective}

\sense{v.}{正确客观地看待}
\ex{Try to keep the problem in perspective — it's not as serious as it seems.}
\col{put sth in(to) perspective 全面看待}
\ex{Travelling puts your daily worries into perspective.}

\begin{origin}
perspective 来自拉丁语 \textit{perspicere}（看透）：\textit{per-}（穿过）+ \textit{specere}（看），即绘画中的"透视法"。"放进透视里看"就是把事物放到正确的位置上看——全局观。
\end{origin}

\vw{celebrity}{/səˈlebrəti/}{n.}

\sense{n.}{名人}
\col{the world of celebrity 名人圈}
\ex{The hotel is popular with celebrities.}

\sense{n.}{名望}
\ex{He gained celebrity as a young novelist.}

\col{celebrity culture 名人文化}

\begin{origin}
来自拉丁语 \textit{celebritas}（众多、热闹），源于 \textit{celeber}（人多的、著名的）。同源词：\textbf{celebrate}（庆祝：众人聚集的场合）。"名人"就是"众人皆知、所到之处都热闹的人"。
\end{origin}

\vw{humility}{/hjuːˈmɪləti/}{n.}

\sense{n.}{谦逊，谦虚}
\col{accept sth with humility 谦逊地接受}
\ex{She accepted the award with humility.}
\col{humble (adj.) 谦卑的}
\ex{Despite his success, he remained humble.}

\begin{origin}
来自拉丁语 \textit{humilitas}（谦卑），源于 \textit{humilis}（低下的），词根 \textit{humus}（泥土、地面）。同源词：\textbf{humble}、\textbf{humiliate}（羞辱：把人按到土里）、\textbf{exhume}（掘出）。"低到贴近地面"就是谦卑——与"趾高气扬"（look down one's nose）正好相反。
\end{origin}

\vw{grace}{/ɡreɪs/}{n.}

\sense{n.}{高雅，文雅；优美}
\col{with grace 优雅地}
\ex{She danced with grace and ease.}

\sense{n.}{风度，体面}
\col{have the grace to do sth 有风度地做某事}
\ex{He had the grace to apologize.}

\col{a grace period 宽限期}
\ex{You have a two-week grace period to pay the bill.}
\col{by the grace of God 承蒙上帝恩典}

\begin{origin}
来自拉丁语 \textit{gratia}（恩惠、喜爱），源于 \textit{gratus}（令人愉悦的、感激的）。同源词：\textbf{grateful}（感激的）、\textbf{gratitude}（感恩）、\textbf{congratulate}（祝贺：共享喜悦）、\textbf{gratis}（免费的）。"令人愉悦的品质"——优雅、风度与恩典同源。
\end{origin}

\vw{dignity}{/ˈdɪɡnəti/}{n.}

\sense{n.}{尊严；自尊；高贵}
\col{with dignity 有尊严地}
\ex{She accepted her defeat with dignity.}
\col{human dignity 人的尊严}
\ex{Everyone has the right to live with dignity.}
\col{beneath one's dignity 有失身份}
\ex{He considered the job beneath his dignity.}

\begin{origin}
来自拉丁语 \textit{dignitas}（价值、尊严），源于 \textit{dignus}（值得的）。同源词：\textbf{dignify}（使有尊严）、\textbf{indignant}（愤慨的：in- 不 + dignus $\to$ 觉得此事不配、令人不齿 $\to$ 义愤）、\textbf{deign}（屈尊）。"值得尊重"——尊严。
\end{origin}

\vp{find one's way into}

\sense{v.}{偶然来到，无意中处于}
\ex{How did these documents find their way into the wrong hands?}
\col{find one's way to 设法到达}
\ex{We eventually found our way to the hotel.}
\col{find its way into（水、光线等）自然进入}
\ex{Water found its way through the cracks in the roof.}

\begin{origin}
find one's way 字面义"找到自己的路"；主语可以是无生命事物（文件、水、光线），表"不知不觉地到达"，是英语中"无生命主语 + 反身指代"的典型用法。
\end{origin}

\vw{league}{/liːɡ/}{n.}

\sense{n.}{联赛}
\col{the Premier League（英）超级联赛}
\ex{Their team was relegated from the league last season.}

\sense{n.}{等级，级别}
\col{in a league of one's own 独领风骚}
\ex{As a striker, he is in a league of his own.}
\col{out of one's league 高不可攀（口语）}
\ex{She's out of my league.}

\col{a league table 联赛积分榜}

\begin{origin}
来自拉丁语 \textit{ligare}（捆绑）$\to$ 意大利语 \textit{lega}（联盟）。同源词：\textbf{ligament}（韧带：把骨头绑在一起）、\textbf{oblige}（ob- 绑住 $\to$ 迫使、使感激）、\textbf{religion}（re- 再 + ligare $\to$ 重新联结神与人 $\to$ 宗教）。"绑在一起的队伍"就是联赛。
\end{origin}

\vw{opposing}{/əˈpəʊzɪŋ/}{adj.}

\sense{adj.}{对抗的；相反的}
\col{the opposing team/side 对方球队}
\ex{Our strikers broke through the opposing defence.}
\col{opposing views/opinions 相反的观点}
\ex{The two candidates hold opposing views on taxation.}

\begin{origin}
oppose 来自拉丁语 \textit{opponere}（置于对面）：\textit{ob-}（对着）+ \textit{ponere}（放置）。同源词：\textbf{opponent}（对手）、\textbf{position}（位置）、\textbf{compose}（com- 一起 + poser 放置 $\to$ 组合）。"摆在对面的人或观点"就是对抗的。
\end{origin}

\vw{net}{/net/}{n.}

\sense{n.}{(the net) 球门网}
\ex{He headed the ball into the net.}
\col{find the back of the net 进球（习语）}

\sense{n.}{网，网状物}
\col{a fishing net 渔网}
\ex{They caught the fish in a large net.}

\col{the Net 互联网（拓展）}
\ex{More and more people get news from the Net.}

\begin{origin}
源自古英语 \textit{net(t)}（网），日耳曼语（同源德语 \textit{Netz}）。20 世纪由"网状物"（network 网络）引申指互联网——the Net。
\end{origin}

\vp{move the goalposts}

\sense{v.}{改变条件（或规则）}
\ex{They moved the goalposts again, changing the rules halfway through the game.}
\ex{The landlord keeps moving the goalposts on the rent.}

\begin{origin}
来自足球：goalposts 是球门柱。趁比赛进行时偷移球门柱，比赛自然没法打——引申为"过程中单方面改变规则、出尔反尔"。
\end{origin}

\vw{insurance}{/ɪnˈʃʊərəns; ɪnˈʃɔːrəns/}{n.}

\sense{n.}{保险；保险业；保险费}
\col{health/life/car insurance}
\ex{He took out life insurance last year.}
\col{an insurance policy 保单}
\ex{Read the terms of the insurance policy carefully.}
\col{insurance against sth 防范……的保障}
\ex{Keep a backup drive as insurance against data loss.}

\begin{origin}
insure 来自古法语 \textit{enseurer}（确保），源于 \textit{seur}（安全的），再追溯到拉丁语 \textit{securus}——即 \textbf{secure}（secure：se- 无 + cura 忧虑 $\to$ 无忧的）。同源词：\textbf{sure}、\textbf{secure}、\textbf{ensure}、\textbf{assure}。"让人安心"——保险的本质。
\end{origin}

\vw{complain}{/kəmˈpleɪn/}{vi. \& vt.}

\sense{vi. \& vt.}{抱怨，投诉}
\col{complain about sth}
\ex{She complained about the noise from the neighbours.}
\col{complain to sb}
\ex{We complained to the manager about the service.}
\col{complain of sth（病人）主诉（病痛）}
\ex{He complained of chest pains.}

\col{complaint (n.) 抱怨；投诉}
\ex{She filed a formal complaint.}

\begin{origin}
来自古法语 \textit{complaindre}（哀叹），源于拉丁语 \textit{complangere}（捶胸痛哭）：\textit{com-}（一起）+ \textit{plangere}（捶打、悲叹）。同源词：\textbf{plaintive}（哀怨的）、\textbf{plaintiff}（原告：plaint 起诉 + -iff $\to$ 提起诉讼的人）。"抱怨"最初是捶胸顿足的哭诉。
\end{origin}

\vp{score an own goal}

\sense{v.}{进乌龙球；无意中做让自己吃亏的事}
\ex{The defender scored an own goal in the last minute.}
\ex{By cutting training budgets, the club scored an own goal.}

\begin{origin}
own goal 是足球术语，指把球踢进自家球门（own 自己的）。引申为"搬起石头砸自己的脚"——本意帮自己，结果害自己。
\end{origin}

\vw{unintentionally}{/ˌʌnɪnˈtenʃənəli/}{adv.}

\sense{adv.}{无意地，非故意地}
\ex{He unintentionally revealed the secret.}
\col{反义：intentionally / on purpose 故意地}
\ex{She accidentally, not intentionally, broke the vase.}

\begin{origin}
来自拉丁语 \textit{intendere}（伸向、意欲）：\textit{in-}（向）+ \textit{tendere}（伸）$\to$ intention（意图，把心思伸向某处）。un-（不）+ intentional + -ly。同源词：\textbf{intend}、\textbf{tension}（tendere 拉伸）。
\end{origin}

\vw{council}{/ˈkaʊnsl/}{n.}

\sense{n.}{政务委员会，地方议会；市政服务机构；委员会}
\col{the city council 市议会}
\ex{The city council approved the new housing plan.}
\col{the UN Security Council 联合国安理会}
\col{a council meeting}
\ex{The council meets once a month.}

\begin{lex}{council / counsel}
\textbf{council} 是"委员会、议会"（\textit{city council}）；\textbf{counsel} 是"律师；忠告"（\textit{counsel for the defence} 辩护律师；\textit{advice and counsel}）。拼写仅差一个字母，常被混用。
\end{lex}

\begin{origin}
来自拉丁语 \textit{concilium}（集会）：\textit{com-}（一起）+ \textit{calare}（大声召集）——与 \textbf{calendar}（历法）同源（calends 古罗马的朔日，须"宣布"的日子）。注意 council 与 counsel 词源不同：counsel 来自 \textit{consilium}（商议）。
\end{origin}

\vw{backfire}{/ˌbækˈfaɪə(r)/}{vi.}

\sense{vi.}{产生事与愿违的不良（或危险）后果}
\col{backfire on sb}
\ex{The plan backfired on him badly.}
\ex{Their strategy to raise prices backfired when sales collapsed.}

\begin{origin}
由 back（回）+ fire（点火）构成，本指内燃机汽缸"逆火"——点火反向爆炸。借喻"计划像引擎一样反着来"，结果适得其反。
\end{origin}

\vw{ballpark}{/ˈbɔːlpɑːk/}{n.}

\sense{n.}{棒球场}
\ex{They watched the game at the local ballpark.}

\sense{n.}{变动范围}
\col{a ballpark figure/estimate 大致相近的数字/估计}
\ex{Can you give me a ballpark figure for the cost?}
\col{in the ballpark 差不多，大致正确}
\ex{His estimate is in the ballpark.}
\ex{The two offers are in the same ballpark.}

\begin{origin}
ball（球）+ park（场地）——棒球场。棒球场大小都差不多，"在同一球场上"即"数量级相近"，遂引申出"大致范围"：a ballpark estimate（大致相近的估计）。
\end{origin}

\vw{venue}{/ˈvenjuː/}{n.}

\sense{n.}{活动场地}
\col{a venue for sth}
\ex{The stadium is the venue for the opening ceremony.}
\col{a concert/wedding venue 音乐会场地；婚礼场地}
\col{change the venue 更换场地}

\begin{origin}
来自古法语 \textit{venue}（到来），源于拉丁语 \textit{venire}（来）。同源词：\textbf{advent}（到来）、\textbf{event}（事件）、\textbf{avenue}（林荫大道：人们走来的路）。venue 就是"人们赶来聚集的地方"。
\end{origin}

\vp{throw sb a curveball}

\sense{v.}{给某人出难题}
\ex{The interviewer threw me a curveball with that question.}

\sense{n.}{curveball 曲线球（棒球术语）}
\ex{The pitcher's curveball broke sharply to the left.}

\begin{origin}
棒球术语：curveball（曲线球）投出后球路突然弯曲，让击球手措手不及。引申为"突然而至的难题、意料之外的情况"。
\end{origin}

\vw{handle}{/ˈhændl/}{vt. \& n.}

\sense{vt.}{处理，应付}
\col{handle a problem/situation}
\ex{She handled the crisis brilliantly.}
\col{handle sth 负责，经营}
\ex{He handles the company's overseas business.}

\sense{vt.}{控制，操纵}
\col{handle a machine}
\ex{This car handles well on wet roads.}

\sense{n.}{把手，拉手；柄}
\col{a door handle}
\ex{Turn the handle to open the door.}
\ex{the handle of a knife}

\begin{origin}
源自古英语 \textit{hand}（手）+ 指小后缀 -le——"手上的小物件"就是把手。同源词：\textbf{hand}、\textbf{handy}（手边的、灵巧的）。"用把手操控"转出动词义"处理、应付"。
\end{origin}

\vp{three strikes and you are out}

\sense{v.}{三振出局}
\ex{In baseball, three strikes and you are out.}
\col{引申：三次犯错即被淘汰}
\ex{Under the new policy, three strikes and you are out — a third offence means dismissal.}

\begin{origin}
棒球规则：击球手三次挥空（strike）即被判出局（out）。此规则被借入司法领域，指"三犯出局"的严厉处罚政策（three-strikes law）。
\end{origin}

\vp{below the belt}

\sense{adj. \& adv.}{不公正的，伤人的}
\ex{That remark was below the belt.}
\ex{Hitting below the belt is against the rules of boxing.}

\begin{origin}
拳击规则：腰带（belt）以下为禁击部位，击打腰部以下即犯规。"击中腰带以下"引申为"不公正、阴险伤人"。
\end{origin}

\vw{remark}{/rɪˈmɑːk/}{n. \& vi. \& vt.}

\sense{n.}{谈论，言论}
\col{make a remark about sth}
\ex{He made a rude remark about her clothes.}
\col{a casual/offhand remark 随口之言}

\sense{vt.}{说起，评论}
\col{remark that ...}
\ex{She remarked that the room was too cold.}
\col{remark on/upon sth 就……发表评论}
\ex{He remarked on the quality of the food.}

\begin{origin}
来自法语 \textit{remarquer}（做记号、注意）：\textit{re-}（再）+ \textit{marquer}（做记号），后者源自日耳曼语"标界"。同源词：\textbf{mark}（记号）。"再做个记号"→"注意到"→"把注意到的事说出来"就是评论。
\end{origin}

\vp{throw in the towel}

\sense{v.}{认输，承认失败}
\ex{After months of losses, the company finally threw in the towel.}
\ex{He refused to throw in the towel despite three failed attempts.}

\begin{origin}
拳击术语：教练见选手不敌，把毛巾（towel）扔进拳台表示认输。引申为"放弃、认栽"。同义表达：give up、throw in the sponge（更早的说法）。
\end{origin}

\vw{literally}{/ˈlɪtərəli/}{adv.}

\sense{adv.}{按字面，字面上}
\col{translate literally 直译}
\ex{The name literally means "light-bringer".}
\col{反义：figuratively 比喻地}

\sense{adv.}{真正地，确实地}
\ex{The city was literally destroyed by the earthquake.}

\begin{origin}
来自拉丁语 \textit{littera}（字母）$\to$ literal（照字母的）$\to$ literally。同源词：\textbf{letter}（字母、信）、\textbf{literature}（文学）、\textbf{literate}（识字的）。"照字母一字不差"就是按字面。
\end{origin}

\begin{gram}[literally 的口语滥用]
现代口语中 literally 常被用作无实义强调词：\textit{I literally died laughing.}（我简直笑死了）——字面义与事实矛盾（人没真死）。权威词典（如《韦氏》）收录了这一用法，但写作时应避免，以免产生歧义或显得夸张失当。
\end{gram}

\vw{iron}{/ˈaɪən/}{n. \& vt. \& vi.}

\sense{n.}{铁}
\ex{Iron is a strong and cheap metal.}
\col{iron (adj.) 铁制的；坚强的（iron will 钢铁意志）}

\sense{n.}{熨斗}
\col{an electric iron 电熨斗}

\sense{vt. \& vi.}{（用熨斗）熨，烫平}
\ex{She ironed his shirt before the interview.}
\col{iron sth out 消除（分歧、困难等）}
\ex{We need to iron out the differences between the two plans.}

\begin{gram}[物质名词的具体化]
\textit{iron}（铁）是不可数物质名词；表示"一件铁器"时变为可数：\textit{an iron}（熨斗）。同类：\textit{glass}（玻璃）→ \textit{a glass}（玻璃杯）；\textit{paper}（纸）→ \textit{a paper}（报纸）；\textit{room}（空间）→ \textit{a room}（房间）。物质名词"装上冠词"即指"一件由该物质制成的东西"。
\end{gram}

\begin{origin}
源自古英语 \textit{isern/iren}，日耳曼语（同源德语 \textit{Eisen}）。词源可追溯到原始印欧语，一说与"天上的金属"有关——早期人类用陨铁打制兵器，许多语言的"铁"都写作"天降之物"（如希腊语 \textit{sideros} 关联"星辰"）。
\end{origin}

% ---------------- 选收的未标 * 词条 ----------------

\vw{compete}{/kəmˈpiːt/}{vi.}

\sense{vi.}{参加比赛；竞争}
\col{compete with/against sb}
\ex{The young athlete competed against the best runners in the world.}
\col{compete for sth}
\ex{More than 200 companies competed for the contract.}
\col{compete in a race/tournament}
\ex{She will compete in the 100-metre race.}

\col{competition (n.) 比赛；竞争}

\begin{origin}
来自拉丁语 \textit{competere}（一同追求）：\textit{com-}（一起）+ \textit{petere}（追求）。同源词：\textbf{petition}（请愿：petere 请求）、\textbf{competent}（competere "共同适配"$\to$ 能胜任的）。"一起追同一个目标"就是竞争。
\end{origin}

\vw{motto}{/ˈmɒtəʊ/}{n.}

\sense{n.}{(pl. mottoes, mottos) 格言，座右铭}
\col{a school motto 校训}
\ex{The Olympic motto is "Faster, Higher, Stronger".}
\ex{His motto is "Never give up".}

\begin{origin}
来自意大利语 \textit{motto}（话语），源于拉丁语 \textit{muttum}（咕哝），是拟声词——同源词 \textbf{mutter}（嘟囔）。从"咕哝的话语"到"座右铭"，一词从低声细语升格为格言。
\end{origin}

\vw{boundary}{/ˈbaʊndri/}{n.}

\sense{n.}{边界，分界线}
\col{the boundary between A and B}
\ex{The river forms the boundary between the two countries.}
\col{push the boundaries 突破界限}
\ex{The artist pushed the boundaries of traditional painting.}

\begin{origin}
bound（界限）来自古法语 \textit{bodne}（界标、疆界），与动词 \textbf{bind}（捆绑）拼写相近但词源无关。同源词：\textbf{bound}（界限）、\textbf{boundary}。"立界标划定的线"就是边界。
\end{origin}

\vw{marathon}{/ˈmærəθən/}{n.}

\sense{n.}{马拉松赛跑；马拉松式的活动（或工作）}
\col{run a marathon}
\ex{He ran his first marathon in under four hours.}
\col{a marathon session/meeting 马拉松式的会议}
\ex{After a marathon session, the committee reached a decision.}

\begin{origin}
来自希腊地名 \textit{Marathōn}（马拉松平原）。公元前 490 年雅典人在马拉松击败波斯人，传令兵费迪皮迪兹跑回雅典报捷后力竭身亡。1896 年首届现代奥运会设马拉松项目纪念此事，距离约 42.195 公里。"马拉松"由此成为"长跑"乃至"持久战"的代名词。
\end{origin}

\vw{opponent}{/əˈpəʊnənt/}{n.}

\sense{n.}{对手，竞争者；反对者}
\col{defeat one's opponent}
\ex{She defeated her opponent in the final.}
\col{a political opponent 政敌}
\col{opponents of the plan 计划的反对者}

\begin{origin}
来自拉丁语 \textit{opponens}（置于对面的人），源于 \textit{opponere}：\textit{ob-}（对着）+ \textit{ponere}（放置）。同源词：\textbf{oppose}（反对）、\textbf{opposing}。"站在对面的人"就是对手。
\end{origin}

\vw{idiom}{/ˈɪdiəm/}{n.}

\sense{n.}{习语，成语}
\col{a common idiom}
\ex{"Kick the bucket" is a colourful idiom meaning "die".}
\col{idiomatic (adj.) 地道的}
\ex{an idiomatic expression}

\begin{origin}
来自希腊语 \textit{idiōma}（特有之物），源于 \textit{idios}（自己的、特有的）。同源词：\textbf{idiot}（idios：只顾自己、不与人交往的人 $\to$ 蠢人）、\textbf{idiosyncrasy}（个人癖性）。习语就是"某种语言特有的说法"。
\end{origin}

\end{document}
